% !TEX root = ../prise-en-main-medical-vrp.tex
\section{Démarrage en 5 minutes}
\label{sec:demarrage}

\subsection{Connexion}

\begin{workflowstep}[Étape 1 — Accéder à l'application]
  Ouvrez l'URL fournie par votre administrateur. Page publique → \menu{Se connecter}. Saisissez e-mail et mot de passe → vous arrivez sur \texttt{/home} (liste des \textbf{structures}).
\end{workflowstep}

\screenshot[keepaspectratio]{01-home-structures.png}{Accueil — liste des structures avec libellés médicaux}

\subsection{Vérifier le contexte métier}

\begin{workflowstep}[Étape 2 — Confirmer le profil médical]
  Menu \menu{Paramètres} → \menu{Entreprise}. Le champ \textbf{Contexte métier} doit afficher \emph{Médical}. Sinon, demandez à l'administrateur de la modifier (seul un admin peut le faire).
\end{workflowstep}

%L'interface bascule alors sur les locales \texttt{fr\_medical} / \texttt{en\_medical} / \texttt{it\_medical}~: menus \menu{Prestations}, \menu{Patients}, \menu{Pilotage d'activité}, etc.

\subsection{Premier tour de l'interface}

\begin{table}[htbp]
  \centering
  \caption{Navigation essentielle (profil médical)}
  \begin{tabular}{@{}ll@{}}
    \toprule
    \textbf{Menu latéral} & \textbf{Rôle} \\
    \midrule
    \menu{Agenda} & Planifier et modifier les \seance{}s \\
    \menu{Trésorerie} & Factures, banque, rapprochement \\
    \menu{Pilotage d'activité} & Liste des \structure{}s, KPIs \\
    \bottomrule
  \end{tabular}
\end{table}

Onglets du module pilotage~: \menu{Pilotage d'activité} · \menu{Structures} · \menu{Prestations} · \menu{Patients}.

\screenshot[keepaspectratio]{02-navigation.png}{Barre latérale et onglets — terminologie médicale}

\subsection{Mode Édition}

Activez le mode \textbf{Édition} dans le bandeau supérieur pour créer ou modifier des données. En \textbf{Consultation}, l'application est en lecture seule.
